\appendix
\section{支撑材料与复现清单}
\footnotesize\sloppy
\begin{longtable}{p{0.31\linewidth}p{0.55\linewidth}}
\caption{当前版本的可复现文件}\label{tab:files}\\
\toprule 文件 & 用途 \\
\midrule\endfirsthead
\toprule 文件 & 用途 \\
\midrule\endhead
\nolinkurl{src/cumcm_b/geometry.py} & 楔形交、凸包、直径、最小包围圆和光学方格 \\
\nolinkurl{src/cumcm_b/policy.py} & q3/q4 覆盖策略、主动测向、联合巡回 \\
\nolinkurl{src/cumcm_b/protocol.py} & 串行接口、计时、幂等请求和失败记录 \\
\nolinkurl{scripts/benchmark_q3_strategies.py} & q3 共同种子配对实验 \\
\nolinkurl{scripts/benchmark_q4_optimization.py} & q4 三种巡回策略配对实验 \\
\nolinkurl{results/synthetic/} & 本地合成实验逐局结果和汇总 \\
\nolinkurl{results/practice/summary.csv} & 匿名官方演练汇总 \\
\nolinkurl{docs/} & 题面转译、证明、边界和复现说明 \\
\bottomrule
\end{longtable}

\section{核心算法摘录}
\begin{lstlisting}[language=Python,caption={安全清除判定的核心逻辑}]
def should_clear(belief):
    polygon = belief.outer_polygon()
    center, radius = enclosing_circle(polygon)
    return radius <= 19.75, center

for channel in pending_channels:
    feedback = measure(channel, current_position)
    update_belief(channel, feedback)
    if feedback.kind == "near":
        clear(channel, current_position)
    elif should_clear(belief[channel])[0]:
        clear(channel, should_clear(belief[channel])[1])
    elif active_queries[channel] >= 6:
        optical_fallback(channel, belief[channel])
\end{lstlisting}
完整源码、逐动作日志和正式测试原始文件不嵌入论文，以仓库当前提交和官方要求的支撑材料为准；提交前重新核对版本号、哈希和匿名化范围。
